% TOP_PROBLEM: {"id": 1955, "title": "Erdős Problem #145", "category_id": 2, "category": {"id": 2, "name": "combinatorics", "display_name": "Combinatorics", "description": "Counting problems, graph theory, discrete structures.", "slug": "combinatorics", "order_index": 2, "created_at": "2026-07-31T15:26:25.671Z"}, "status": "open"}
% FIRST_POSTED: null
% ENTRY_KIND: statement_only
% Imported from: unsolvedmath_additional_500.tex; UnsolvedMath ID 1955
% !TeX program = lualatex
% Compile twice: lualatex 1955.tex
% Self-contained: no external bibliography, images, or \input files.
% Numeric section labels and visible numbers are original UnsolvedMath IDs.
\documentclass[11pt,a4paper]{article}
\usepackage[margin=24mm,headheight=14pt]{geometry}
\usepackage{fontspec}
\setmainfont{Latin Modern Roman}
\setsansfont{Latin Modern Sans}
\setmonofont{Latin Modern Mono}
\usepackage{amsmath,amssymb,mathtools}
\usepackage{unicode-math}
\setmathfont{Latin Modern Math}
\usepackage{microtype}
\usepackage{enumitem,xurl,needspace}
\usepackage[dvipsnames]{xcolor}
\usepackage{hyperref}
\hypersetup{unicode=true,colorlinks=true,linkcolor=MidnightBlue,urlcolor=MidnightBlue,
 pdftitle={UnsolvedMath Problem 1955},pdfauthor={Source-annotated compilation},bookmarksnumbered=true}
\usepackage{bookmark,tocloft,titlesec,fancyhdr}
\setlength{\cftsecnumwidth}{4.2em}
\cftsetpnumwidth{2.5em}
\cftsetrmarg{4em}
\titleformat{\section}{\Large\bfseries\raggedright}{\thesection}{0.7em}{}
\pagestyle{fancy}\fancyhf{}
\fancyhead[L]{\small\sffamily Additional UnsolvedMath statements}
\fancyhead[R]{\small Original catalogue IDs}
\fancyfoot[C]{\thepage}
\setlength{\parindent}{0pt}
\setlength{\parskip}{5pt plus 1pt}
\setlength{\emergencystretch}{3em}
\setcounter{tocdepth}{1}\setcounter{secnumdepth}{1}
\setlist[itemize]{leftmargin=3em,itemsep=4pt,topsep=4pt,parsep=0pt}
\allowdisplaybreaks
\newcommand{\N}{\mathbb{N}}
\newcommand{\Z}{\mathbb{Z}}
\newcommand{\Q}{\mathbb{Q}}
\newcommand{\R}{\mathbb{R}}
\newcommand{\C}{\mathbb{C}}
\newcommand{\F}{\mathbb{F}}
\newcommand{\A}{\mathbb{A}}
\newcommand{\PP}{\mathbb{P}}
\newcommand{\E}{\mathbb{E}}
\newcommand{\eps}{\varepsilon}
\newcommand{\dd}{\,\mathrm d}
\newcommand{\sref}[2]{\hyperlink{source:#1:#2}{[S#2]}}
\newif\ifshowcatalogue
\showcataloguetrue
\newcommand{\cataloguescope}[1]{\ifshowcatalogue
 \paragraph{Statement recorded in the supplied source.}
 #1\fi}
\begin{document}
% UnsolvedMath ID 1955; source catalogue code EP-145

\setcounter{section}{1954}

\section{Limiting Moments of Squarefree Gaps}\label{1955}

{\small\sffamily UnsolvedMath ID 1955 · Combinatorics}

\subsection{Definitions and mathematical statement}

\paragraph{Shared notation.}Unless a section says otherwise, $\N=\{1,2,\ldots\}$,
$\Z$ is the ring of integers, $\Q,\R,\C$ are the rational, real, and complex
fields, $[N]=\{1,\ldots,\lfloor N\rfloor\}$, and $\log$ is the natural
logarithm. For real functions of a parameter tending to infinity,
$f\ll g$ and $f=O(g)$ mean $|f|\leq Cg$ eventually for some fixed $C$;
$f\gg g$ reverses this comparison; $f=o(g)$ means $f/g\to0$;
$f\sim g$ means $f/g\to1$; and $f\asymp g$ means both order bounds.
A subscript on $O$ or $\ll$ specifies allowed constant dependence.
The expression $n^{o(1)}$ in an upper bound means growth smaller than
$n^\epsilon$ for each fixed $\epsilon>0$. Local definitions govern any
exception, finite-field convention, or use of the same symbol for another
object. An asymptotic claim is not automatically an effective algorithm.



An integer is squarefree if it is divisible by no prime square. The gap attached to $s_n\leq x$ is the full gap $s_{n+1}-s_n$, even when $s_{n+1}>x$. The exponent is fixed before taking the limit. \sref{1955}{1} \sref{1955}{2}

\paragraph{Statement recorded in the supplied source.}
Let $s_1<s_2<\cdots$ be the sequence of squarefree numbers. Is it true that, for any $\alpha \geq 0$, $ \lim_{x\to \infty}\frac{1}{x}\sum_{s_n\leq x}(s_{n+1}-s_n)^\alpha $ exists? \sref{1955}{1}

\subsection{Short English statement}

Do all nonnegative moments of gaps between consecutive squarefree integers have a limit when accumulated per unit interval length?

\subsection{Sources}

{\small

\begin{itemize}

\item[\hypertarget{source:1955:1}{S1}] ULAM.AI, \emph{UnsolvedMath}, supplied \texttt{problems.json}, record 1955 (EP-145), statement and background. The embedded literature-review date is 2026-08-17; that is the archive’s review date, not a fresh certification. Dataset landing page: \url{https://huggingface.co/datasets/ulamai/UnsolvedMath}.

\item[\hypertarget{source:1955:2}{S2}] T. F. Bloom, \emph{Erdős Problems}, Problem 145, \url{https://www.erdosproblems.com/145}. Reference imported from the supplied record; the live problem page was not independently checked for this compilation.

\item[\hypertarget{source:1955:3}{S3}] Tsz Ho Chan, On moments of gaps between consecutive square-free numbers, Moscow Journal of Combinatorics and Number Theory 12 (2023), 287-295. \url{https://doi.org/10.2140/moscow.2023.12.287}. Bibliographic information imported from the supplied record.

\item[\hypertarget{source:1955:4}{S4}] Erd\"{o}s, P., Some problems and results in elementary number theory. Publ. Math. Debrecen (1951), 103-109. Bibliographic information imported from the supplied record.

\item[\hypertarget{source:1955:5}{S5}] Greaves, G. R. H. and Harman, G. and Huxley, M. N., Sieve Methods, Exponential Sums, and their Applications in Number Theory. (1997). Bibliographic information imported from the supplied record.

\item[\hypertarget{source:1955:6}{S6}] Granville, Andrew, {$ABC$} allows us to count squarefrees. Internat. Math. Res. Notices (1998), 991--1009. Bibliographic information imported from the supplied record.

\end{itemize}

}

\label{end:1955}
\end{document}
